\documentclass[11pt,a4paper]{amsart}
\usepackage[a4paper,landscape,margin=20mm]{geometry}
\usepackage{amsmath,amssymb,amsthm,graphicx,xcolor}
\usepackage[no-math]{fontspec}
\setmainfont{lmroman10-regular.otf}[BoldFont=lmroman10-bold.otf,ItalicFont=lmroman10-italic.otf,BoldItalicFont=lmroman10-bolditalic.otf,SmallCapsFont=lmromancaps10-regular.otf]
\usepackage[inline,shortlabels]{enumitem}
\theoremstyle{plain}

\newtheorem{thm}{Theorem}[section]
\newtheorem{cor}[thm]{Corollary}
\newtheorem{lem}[thm]{Lemma}
\newtheorem{slem}[thm]{Sublemma}
\newtheorem{prop}[thm]{Proposition}

% the following environments keep the roman font:
\theoremstyle{definition}

\newtheorem{asm}[thm]{Assumption}
\newtheorem{rem}[thm]{Remark}
\newtheorem{rems}[thm]{Remarks}
\newtheorem{exa}[thm]{Example}
\newtheorem{exas}[thm]{Examples}
\newtheorem{defi}[thm]{Definition}
\newtheorem{conv}[thm]{Convention}
\newtheorem{conj}[thm]{Conjecture}
\newtheorem{prob}[thm]{Problem}
\newtheorem{oprob}[thm]{Open Problem}
\newtheorem{oprobs}[thm]{Open Problems}
\newtheorem{algo}[thm]{Algorithm}
\newtheorem{obs}[thm]{Observation}
\newtheorem{desc}[thm]{Description}
\newtheorem{fact}[thm]{Fact}
\newtheorem{qu}[thm]{Question}
\newtheorem{oqu}[thm]{Open Question}
\newtheorem{pty}[thm]{Property}
\newtheorem{clm}[thm]{Claim}
\newtheorem{nota}[thm]{Notation}
\newtheorem{com}[thm]{Comment}
\newtheorem{coms}[thm]{Comments}

%% goal is to omit parens from note, place cite as part of theorem number
\newtheoremstyle{defC}%
  {6pt}% (space above)
  {6pt}% (space below)
  {\normalfont}% (body font)
  {}% (indent amount)
  {\bfseries}% {theorem head font}
  {{\bfseries .}}% {punctuation after theorem head}
  {5pt plus 1pt minus 1pt}% {space after theorem head}
  {\thmname{#1} \thmnumber{#2} \thmnote{\normalfont#3}}% {theorem head spec}
\theoremstyle{defC}
\newtheorem{defiC}[thm]{Definition}
\newtheorem{remC}[thm]{Remark}
\newtheorem{exaC}[thm]{Example}

\newtheoremstyle{thmC}%
  {6pt}% (space above)
  {6pt}% (space below)
  {\itshape}% (body font)
  {}% (indent amount)
  {\bfseries}% {theorem head font}
  {{\bfseries .}}% {punctuation after theorem head}
  {5pt plus 1pt minus 1pt}% {space after theorem head}
  {\thmname{#1} \thmnumber{#2} \thmnote{\normalfont#3}}% {theorem head spec}

\theoremstyle{thmC}
\newtheorem{thmC}[thm]{Theorem}
\newtheorem{propC}[thm]{Proposition}
\newtheorem{lemC}[thm]{Lemma}

% switch back to standard theorem style for document specific theorems
\theoremstyle{plain}

\numberwithin{equation}{section}


%% Changed applied up-to 21st april 2021
    %% 
\newif\iflong\longtrue
\newif\iflongF\longFtrue %further, last minutes cut
\newif\ifcutLICS\cutLICStrue
\newif\ifcutLMCS\cutLMCSfalse

\newif\ifcutCBNLMCS\cutCBNLMCSfalse  
 % final cuts for the Call-By-Name part, right before the LMCS submission

\newif\ifnotes\notesfalse

\newif\ifcommentreview\commentreviewfalse

  
\usepackage{hyperref}
\usepackage{cite}
\usepackage{xspace}
% \input{ourAMS}   % this is different between DS and GJ
\usepackage{bm}
\usepackage{amssymb,amsmath,stmaryrd,mathtools}
\DeclareMathAlphabet{\mathbbmsl}{U}{bbm}{m}{sl}
\usepackage{amsthm}
\usepackage{xcolor}
\usepackage{mathpartir}
% ORIGINAL INPUT davidemacros.sty
%% Restored, but to be removed
\newcommand{\wb}{\bsim}
\newcommand{\bsim}{\wbPI} % {\ensuremath{\approx}\xspace}
\newcommand{\bsiml}{\ensuremath{\approx_{\L}}\xspace}
\newcommand{\TE}{\mathrel{\approx_{\rm tr}}}       %trace equivalence
\newcommand{\TEs}{\mathrel{\approx_{\rm tr}^{\rm s}}}       %trace equivalence in op LTS

% utility

\newcommand{\wrt}{with respect to\xspace}  % 

\newcommand{\AG}{T}  % pi agents

\newcommand{\okTR}[2]{{\tt ok}(#1,#2)}  %well formed traces
                                                %  #2 \vdash #1
\newcommand{\complete}{complete\xspace}
\newcommand{\completetrace}{complete-trace\xspace}


\newcommand{\adjust}[1]{
\color{blue}
  \marginpar{$\dwa${\bf BL}} \vskip .1cm 
#1
  \marginpar{$\uparrow${\bf BL}}\vskip .1cm 
\color{black}
}

% \newcommand{\adjust}[1]{}

\newcommand{\proposed}[1]{
\color{violet}
  \marginpar{$\Dwa${\bf VI}}\vskip .1cm 
#1
  \marginpar{$\Uparrow${\bf VI}}\vskip .1cm 
\color{black}
}

\ifnotes
\def\finish#1{\vskip.2cm\noindent{\em #1}%
  \marginpar{$\longleftarrow$}\vskip.2cm}
  \else
 \newcommand{\finish}[1]{}
\fi


\def\note#1{\vskip.2cm\noindent{\em #1}%
  \marginpar{
\begin{tabular}[t]{l}
$\leftarrow$ \\
 (Note for us)
\end{tabular} 
 \vskip.2cm}} 

% for appendixes

\newcommand{\mydef}[1]{
\vskip .2cm 
\noindent 
({\bf Definition~\ref{#1}})
}

\newcommand{\myteo}[1]{
\vskip .2cm 
\noindent 
({\bf Theorem~\ref{#1}})
}



% for WBC traces
\newcommand{\trace}{s}
\newcommand{\tr}{s}
\newcommand{\tracerho}{t}
\newcommand{\tracetheta}{u}
\newcommand{\stack}[2]{#1[#2]}
%\newcommand{\addstack}[2]{(#1,#2)}
\newcommand{\stt}{\sigma}
\newcommand{\stti}{\sigma_i}
\newcommand{\sttp}{\sigma'}
\newcommand{\stto}{\sigma_1}
\newcommand{\sttt}{\sigma_2}




%for operational semantics of pi

\newcommand{\pair}[2]{  #1 }      % 
\newcommand{\pairP}[1]{  #1 }      % 
\newcommand{\pairPP}[1]{ #1 }      % 
% \newcommand{\pair}[2]{\langle  #1, #2\rangle }      % 
% \newcommand{\pairP}[1]{\langle  #1, \phi \rangle }      % 
% \newcommand{\pairPP}[1]{\langle  #1, \phi' \rangle }      % 
\newcommand{\equival}{\equiv_\alpha }
\newcommand{\trans}[1]{\mbox{{\tt #1}}}


 \newcommand{\caserule}[1]{(rule (#1))}

\newcommand{\mydots}{ ,\ldots, }
\newcommand{\aalpha}{\mu}

\newcommand{\LL}{L} % for extended call-by-name  terms


%parallel OGS

\newcommand{\AOGS}{\mbox{{\rm A-OGS}}\xspace}      % 
\newcommand{\COGS}{\mbox{{\rm C-OGS}}\xspace}      % 
\newcommand{\WBOGS}{\mbox{{\rm WB-OGS}}\xspace}      % 
\newcommand{\OGS}{\mbox{{\rm OGS}}\xspace}      % 

\newcommand{\Support}{support\xspace}      % 
\newcommand{\Supports}{supports\xspace}      % 


% stacked processes
\newcommand{\sta}[2]{#1 \ltimes #2}      % pi configuration

 % configurations for OGS parallel 

\newcommand{\set}[3]{\{ #1_{#3}, #2_{#3} \}_{#3}}      
\newcommand{\setI}[4]{\{ #1_{#3}, #2_{#3} \}_{#3\in I}}      

\newcommand{\Dom}[1]{{{\tt dom}}(#1)} 

\newcommand{\appendSYMB}{\otimes}

\newcommand{\appendCon}[2]{#1 \otimes #2} 
\newcommand{\appendRel}[1]{#1^{|}} 
\newcommand{\uptoComp}[1]{\mathrel{#1^{|\star}}} 
\newcommand{\uptoCompRed}[1]{\mathrel{\Rightarrowp #1^{|\star} 
                                   \Leftarrowp}} 

%more on config

\newcommand{\okC}[2]{{\tt se}(#1,#2)}  %compatibility      
\newcommand{\LN}[1]{{\tt ln}(#1) \LNsym}  % loc names      
\newcommand{\LNsym}{\ltimes}  % loc names symbol      

\newcommand{\okCON}[2]{{\tt ok}(#1,#2)}  % compatibility between a configuration and a stack



% names

%\def\Names{{\mathcal N}}
\def\fn#1{\mbox{\tt fn}(#1)}         % free names
\def\bn#1{\mbox{\tt bn}(#1)}         % bound names
%\def\en#1{\mbox{\tt en}(#1)}         % extruded names
\def\n#1{\mbox{\tt n}(#1)}           % names

\def\fv#1{\mbox{\tt fv}(#1)}         % free variables in Lambda-calculus 


% processes

\newcommand{\obj}[1]{\langle #1 \rangle}
\def\nil{0}
\def\res#1{{\boldsymbol \nu} #1\:}   % restriction 
\def\resb#1{({\boldsymbol \nu} #1)}   % restriction  with bracke 
\def\rest#1#2{({\boldsymbol \nu} #1 : #2)\:} % typed restriction

\def\out#1#2{\overline{#1}\obj{#2}}
\def\inp#1#2{#1(#2)}
\def\inpw#1#2{#1(#2)}
\def\bout#1#2{\overline{#1}(#2)}   % bound output
\newcommand{\outb}{\bout}           % bound output
\newcommand{\new}{\res}           % restriction

\newcommand{\inpi}[3]{ \inpw{#1_{#2}}{#3} }%input Prefix indexed
\newcommand{\opi}[3]{ \overline{#1}_{#2}\obj{#3} }%Output Prefix indexed
\newcommand{\bopi}[3]{ \overline{#1}_{#2}({#3}) }%Output Prefix indexed

\def\con#1#2{#1 \langle{#2}\rangle}      % constants
\def\abs#1{(#1). }                    % abstractions




\newcommand{\outC}[1]{\overline{#1}}      % CCS  output
\newcommand{\inpC}[1]{#1}                 % CCS  input


\mathcode`\!="4021 % `!' as prefix operator
\mathcode`\|="326A % `|' as relation operator
\mathcode`\.="602E


\def\con#1#2{#1 \langle{#2}\rangle}      % constants
\def\abs#1#2{(#1). #2}                    % abstractions
\def\absT#1#2{(#1:#2)}                    % typed abstractions
\def\app#1#2{#1\obj{#2}}                    % applications


\newcommand{\pr}{{\mathcal P}}

% FOR CONTEXTS

\newcommand{\ctx}[1]{#1^{\rmm{ctx}}  }   % context closure


\newcommand{\contexthole}{ [ \cdot  ] }      %hole of a context
\newcommand{\ct}[1]{ C_{\rmm{ctx}} \brac{#1} }   %filled context
\newcommand{\qct}{ C_{\rmm{ctx}}  }              %unfilled context  
\newcommand{\ctp}[1]{ C' \brac{#1} }   %filled context
\newcommand{\qctp}{ C'  }              %unfilled context  
\newcommand{\brac}[1]{[#1] }   % brackets for the context hole

\newcommand{\Ep}{E'}              % evaluation contexts


% redefined \rightarrowfill, with lower  labels

\def \rightarrowfill{$\m@th\mathord{\smash-}\mkern-6mu%
  \cleaders\hbox{$\mkern-2mu\mathord{\smash-}\mkern-2mu$}\hfill
  \mkern-6mu\mathord\rightarrow$}
\def \mapstofill{$\m@th\mathord{\smash\mapstochar\smash-}\mkern-6mu%
  \cleaders\hbox{$\mkern-2mu\mathord{\smash-}\mkern-2mu$}\hfill
  \mkern-6mu\mathord\rightarrow$}

% transizioni

\newcommand{\rmm}[1]{{\tt{\lowercase{#1}}}}
\newcommand{\rmmm}[1]{{\tt{{#1}}}}

\newcommand{\genarr}[2]{\mathrel{\stackrel{{#1}}{\longrightarrow}_{\rmm{\lowercase{#2}}}}}  
                                    % parametrised strong arrow
\newcommand{\genArr}[2]{\mathrel{\stackrel{{#1}}{\Longrightarrow}_{\rmm{\lowercase{#2}}}}}  %  
                                    % parametrised weak arrow

\newcommand{\genArcap}[2]{\mathrel{\stackrel{{\widehat{#1}}}{\Longrightarrow}_{\rm{\lowercase{#2}}}}}
                                    % parametrised weak transitions with   hat
\newcommand{\genlongrightarrow}[1]{\mathrel{\longrightarrow_{\rmm{\lowercase{#1}}}}}        
\newcommand{\genLongrightarrow}[1]{\mathrel{\Longrightarrow_{\rmm{\lowercase{#1}}}}}        
\newcommand{\genLeftarrow}[1]{\mathrel{\mbox{}_{\rmm{\lowercase{#1}}}\!\!\Leftarrow}}        
\newcommand{\genRightarrow}[1]{\mathrel{\Rightarrow_{\rmm{\lowercase{#1}}}}}



\newcommand{\longrightarrowG}{\genlongrightarrow{G}}
\newcommand{\LongrightarrowG}{\genLongrightarrow{G}}




\newcommand{\arrG}[1]{\genarr{#1}{G}}   %  generic strong arrow
\newcommand{\ArrG}[1]{\genArr{#1}{G}}   %  generic weak arrow
\newcommand{\ArcapG}[1]{\genArcap{#1}{G}}   %  generic weak arrow hat


\newcommand{\arrPI}[1]{\genarr{#1}{\pi}}   %strong arrow in pi
\newcommand{\ArrPI}[1]{\genArr{#1}{\pi}}   % weak pi
\newcommand{\longrightarrowPI}{\genlongrightarrow{\pi}}
\newcommand{\LongrightarrowPI}{\genLongrightarrow{\pi}}

\newcommand{\arrA}[1]{\genarr{#1}{A}}  %  Alternating OGS LTS
\newcommand{\ArrA}[1]{\genArr{#1}{A}}  %  Alternating OGS LTS 
\newcommand{\longrightarrowA}{\genlongrightarrow{A}}
\newcommand{\LongrightarrowA}{\genLongrightarrow{A}}





\newcommand{\arrC}[1]{\genarr{#1}{C}}   %strong arrow in COGS
\newcommand{\ArrC}[1]{\genArr{#1}{C}}   % weak COGS
\newcommand{\Arcapp}[1]{\genArcap{#1}{C}}   % weak COGS with cap
\newcommand{\longrightarrowp}{\genlongrightarrow{C}}
\newcommand{\Longrightarrowp}{\genLongrightarrow{C}}
\newcommand{\Leftarrowp}{\genLeftarrow{C}}
\newcommand{\Rightarrowp}{\genRightarrow{C}}




\newcommand{\arrWB}[1]{\genarr{#1}{WB}}           %   arrow for WB LTS
\newcommand{\ArrWB}[1]{\genArr{#1}{WB}}          %   weak arrow WB LTS



\newcommand{\arrN}[1]{\genarr{#1}{o\pi}}  %  seq LTS  in pi
\newcommand{\ArrN}[1]{\genArr{#1}{o\pi}}  %  seq LTS in pi


\newcommand{\arr}[1]{\mathrel{\stackrel{{#1}}{\longrightarrow}}}  %  strong arrow
\newcommand{\Arr}[1]{\mathrel{\stackrel{{#1}}{\Longrightarrow}}}  %  weak arrow

                                   
\newcommand{\arcap}[1]{\:\stackrel{{\widehat{#1}}}{\ra}\:}
                                    %strong labelled transitions with hat
\newcommand{\Arcap}[1]{\:\stackrel{{\widehat{#1}}}{\Longrightarrow}\:}
                                    %weak labelled transitions with   hat
                                    

% reductions
\newcommand{\ra}{\longrightarrow}        
\newcommand{\Ra}{\Longrightarrow} 

\newcommand{\raN}{\arrN{}}        
\newcommand{\RaN}{\ArrN{}} 


\newcommand{\outputpr}{output-prioritised\xspace}  %  sequential  LTS pi
\newcommand{\opLTS}{opLTS\xspace}  %  sequential  LTS pi


% equivalences and preorders

\newcommand{\genTE}[1]{\mathrel{\bumpeq_{\rmm{\lowercase{#1}}}}}       %parametrised trace equivalence


\newcommand{\genwb}[1]{\mathrel{\approx_{\rmm{\lowercase{#1}}}}}       %parametrised bisimilarity 
\newcommand{\wbG}{\genwb{G}}       % generic  bisimilarity 
\newcommand{\wbPI}{\genwb{\pi}}       %   bisimilarity  in pi
\newcommand{\wbOPI}{\genwb{o\pi}}       %   bisimilarity  in pi, op LTS
\newcommand{\wbC}{\genwb{C}}       % bisimulation in COGS
\newcommand{\wbA}{\genwb{A}}       % bisimulation in Alternating OGS
\newcommand{\wbWB}{\genwb{wb}}       % bisimulation in Well-Bracketing OGS

\newcommand{\genwbDIV}[2]{\mathrel{\approx^{#2}_{\rmm{\lowercase{#1}}}}}
                                 % parametrised bisimilarity  respecting divergence
\newcommand{\wbPIDIV}{\genwbDIV{\pi}{\Uparrow}}       %   bisimilarity in pi resp. div.


\newcommand{\TEG}{\genTE{G}}       % generic trace equivalence
\newcommand{\TEPI}{\genTE{\pi}}       % trace equivalence in pi
\newcommand{\TEN}{\genTE{o\pi}}     %trace equivalence in the opLTS pi

\newcommand{\TEA}{\genTE{A}}       % trace equivalence in Alternating OGS LTS
\newcommand{\TEp}{\genTE{C}}       %trace equivalence in COGS
\newcommand{\TEwb}{\genTE{wb}}       %trace equivalence in WBOGS

\newcommand{\LR}{\genTE{C\pi}}       %trace equivalence under linear receptiveess



\newcommand{\genTEcomp}[1]{\mathrel{\bumpeq^{\rm c}_{\rm{\lowercase{#1}}}}} %parametrised complete trace equivalence
\newcommand{\genTEcompST}[2]{\mathrel{\bumpeq^{\rm c[#1]}_  {\rm{\lowercase{#2}}}}} 
                 %  parametrised complete trace equivalence under a stack #1 

\newcommand{\TEwbc}{\genTEcomp{WB}}       % complete trace equivalence in WB OGS
% \newcommand{\TEwbc}{\genTEcomp{WBinit}}       % complete trace equivalence in WB OGS on initial conf
\newcommand{\TEc}{\genTEcomp{A}}        % complete traces in AOGS
\newcommand{\TEpc}{\genTEcomp{C}}        % complete traces in COGS LTS
\newcommand{\TEwbcN}[1]{\genTEcompST{#1}{WB}}    % complete traces in WB LTS wrt a stack #1
\newcommand{\TEcN}[1]{\genTEcompST{#1}{A}}       % complete traces wrt #1
\newcommand{\TEpcN}[1]{\genTEcompST{#1}{C}}       %  complete traces wrt #1 in COGS


% expansion
\newcommand{\contr}{\mathrel{\mbox{}_{\pi}\!\!\gtrsim}}
\newcommand{\contrG}{\contr}
\newcommand{\expa}{\mathrel{\lesssim_{\pi}}}
\newcommand{\exn}{\contr} %_{\mathrm{exp}}}
\newcommand{\lexn}{\expa} %_{\mathrm{exp}}}

\newcommand{\expaOP}{\mathrel{\lesssim_{\rmm{o}\pi}}} % expansion in the opLTS 
\newcommand{\contrOP}{\mathrel{\mbox{}_{\rmm{o}\pi}\!\!\gtrsim}}


\newcommand{\interl}[2]{{\tt inter}( #1, #2)}
\newcommand{\interlA}[2]{{\tt inter_{a}}( #1, #2)}
\newcommand{\interlWB}[2]{{\tt inter_{wb}}( #1, #2)}
\newcommand{\interlWBs}[3]{{\tt inter_{wb}^{#1}}( #2, #3)}
\newcommand{\commu}[2]{{\tt comm}( #1, #2)}


%  special relations

\def\R{{\mathcal R}}
\def\RR{\mathrel{\mathcal R}}
\def\S{{\mathcal S}}
\def\SS{\mathrel{\mathcal S}}


% defuguale

\def\defi{\stackrel{\mbox{\scriptsize {\rm{def}}}}{=}} % definition equal
% \newcommand{\Defi}{\mathrel{\triangleq}}           %
%\def\Defi{\stackrel{\mbox{\scriptsize $\triangle$}}{=}} % definition equal


%convergence

\newcommand{\Dwa}{\Downarrow}           % plain convergence
\newcommand{\dwa}{\downarrow}           % plain convergence


\newcommand{\DwaLR}[1]{\Downarrow^{#1}_{\rm{LR}}}          
\newcommand{\DwaN}[1]{\Downarrow^{#1}}           % 
\newcommand{\dwaLR}[1]{\downarrow^{#1}_{\rm{LR}}}  
\newcommand{\dwaN}[1]{\downarrow^{#1}}           



%SPECIAL SYMBOLS

\ifcutLMCS
\def\Mybar{{\center \rule{\hsize}{0.5mm}}} % bar in tables
\else
\def\Mybar{}
\fi
\def\dottedbar{{\center \rule{\hsize}{0.2mm}}} % bar in tables

\def\midd{\; \; \mbox{\Large{$\mid$}}\;\;} % separation symbol in
                                           % grammars               
\def\st{\; : \;}

\newcommand{\til}[1]{\widetilde{#1}}
\newcommand{\tilb}{{\til b}}
\newcommand{\tila}{{\til a}}
\newcommand{\tilc}{\tilde{c}}
\newcommand{\tilx}{\tilde{x}}
\newcommand{\tily}{\tilde{y}}
\newcommand{\tilz}{\tilde{z}}
\newcommand{\tilw}{\tilde{w}}

\newcommand{\veca}{\vec{a}}
\newcommand{\vecb}{\vec{b}}
\newcommand{\vecc}{\vec{c}}

\newcommand{\vecq}{\vec{q}}
\newcommand{\vecp}{\vec{p}}
\newcommand{\vecqi}{\vec{q_i}}
\newcommand{\vecqo}{\vec{q_1}}
\newcommand{\vecqt}{\vec{q_2}}
\newcommand{\vecqp}{\vec{q'}}

\newcommand{\Ki}{\FF_i}
\newcommand{\Ko}{\FF_1}
\newcommand{\Kt}{\FF_2}
\newcommand{\Hi}{\GG_i}
\newcommand{\Ho}{\GG_1}
\newcommand{\Ht}{\GG_2}

\newcommand{\HH}{H}
\newcommand{\KK}{K}
\newcommand{\FF}{F}
\newcommand{\GG}{G}
\newcommand{\DD}{D}
\newcommand{\PP}{P}
\newcommand{\QQ}{Q}

\newcommand{\FFc}{\FF}%{{\mathcal{\FF}}} % configurations in COGS
\newcommand{\GGc}{\GG}% {{\mathcal{\GG}}} % configurations in COGS


\newcommand{\FFb}{{\bm{\FF}}} % configurations with stack
\newcommand{\GGb}{{\bm{\GG}}} % configurations with stack
%\newcommand{\FFb}{{\textbf{\it{\FF}}}} % configurations with stack
%\newcommand{\GGb}{{\textbf{\emph{\GG}}}}


\newcommand{\es}{\eraseStack}

 \newcommand{\cogs}[1]{#1} %{\{#1\}_{\rmm{C}}}
 \newcommand{\ogs}[1]{#1} %{\{#1\}_{\rmm{A}}}
% \newcommand{\cogs}[1]{\uparrow_{\rmm{C}}\!#1}
% \newcommand{\ogs}[1]{\uparrow_{\rmm{A}}\!#1}


% for encoding 

% variant with linear continuations
\newcommand{\encoW}[1]{\cbvSymbNoLin [\! [ #1 ] \! ]  }
\newcommand{\encoWa}[2]{\cbvSymbNoLin \app{[\! [ #1 ] \! ]}{#2}  }
\newcommand{\cbvSymbNoLin}{{\mathcal {W}}}


\newcommand{\enco}[2]{ \app{[\! [ #1 ] \! ]}{#2}  }
\newcommand{\encoO}[1]{{\mathcal {A}} [\! [ #1 ] \! ]}  
\newcommand{\encoOa}[2]{{\mathcal {A}} \app{[\! [ #1 ] \! ]}{#2}  }
\newcommand{\encoI}[1]{\cbvSymb [\! [ #1 ] \! ]  }
\newcommand{\encoIa}[2]{\cbvSymb \app{[\! [ #1 ] \! ]}{#2}  }
\newcommand{\encoIV}[1]{\cbvSymb^\variant [\! [ #1 ] \! ]  }
\newcommand{\encoIVa}[2]{\cbvSymb^\variant \app{[\! [ #1 ] \! ]}{#2}  }
\newcommand{\cbvSymb}{{\mathcal {V}}}


%% Call-By-Name 
\newcommand{\encoN}[1]{{\mathcal {N}} [\! [ #1 ] \! ]  }
\newcommand{\encoNa}[2]{{\mathcal {N}} \app{[\! [ #1 ] \! ]}{#2}  }

\newcommand{\encoS}[1]{ {[\! [ #1 ] \! ]}  } % call-by-name abbreviated
\newcommand{\encoSa}{\enco}                    % call-by-name abbreviated      


\newcommand{\encoSN}{\encoN} 
% used to be: encoS, changed into \encoSN so to avoid abbreviation before LMCS submission 
                 


\newcommand{\encoSP}[2]{ \pairP{\app{[\! [ #1 ] \! ]}{#2}}  } % call-by-name with configurations
\newcommand{\encoSPN}[2]{ \pairP{\app{ \encoSN{ #1}}{#2}}  } 
  % call-by-name with configurations : 
   % changed from encoSP so to add N before square brackets, before LMCS submission

\newcommand{\encP}[1]{ \pairP{{[\! [ #1 ] \! ]}}}   % call-by-name with configurations
\newcommand{\encPP}[2]{ \pair{{[\! [ #1 ] \! ]}}{#2}  } % call-by-name with configurations

\newcommand{\encPPN}[2]{ \pair{{\encoSN{ #1 }}}{#2}  } 
  % call-by-name with configurations : 
   % changed from encoSP so to add N before square brackets, before LMCS submission


%% Call-By-Value 

\newcommand{\enca}[1]{\cbvSymb[\![#1]\!]}

\newcommand{\encoEnvI}[1]{ \cbvSymb[\! [ #1 ] \! ]  } % enc of environments
\newcommand{\encoConI}{\encoEnvI} % enc of config.


\newcommand{\encoEnvV}[1]{{\mathcal {O}} [\! [ #1 ] \! ]  } % enc of environments
\newcommand{\encoConV}{\encoEnvV} % enc of config.


\newcommand{\bind}[1]{({#1})~}

%The optimized encoding (for the soundness proof)
\newcommand{\encoV}[1]{{\mathcal O}[\![#1]\!]}  %\encba
\newcommand{\qencV}{{\mathcal O}}  %\qencba
\newcommand{\encoVa}[2]{\app{\encoV{#1} }{#2}} %\encb

%\newcommand{\encla}[1]{\llbracket #1\rrbracket^{\L}}

% value part of the optimised encoding
\newcommand{\encoVVa}[2]{\app{\encoVV{#1}}{#2}}  %\encbv
\newcommand{\encoVV}[1]{{\mathcal O}^\variant[\![#1]\!]} %\encbva
\newcommand{\variant}{*} %{\mathrm{V}}
%\newcommand{\param}[1]{\langle #1 \rangle}


\newcommand{\equaDS}[4]{ #2& #3& #4&  \mbox{{$#1 $}}}

\newcommand{\nameDS}[1]{{\rn{#1}}}
\newcommand{\rn}[1]{%
  \ifmmode 
    \mathchoice
      {\mbox{\sc #1}}
      {\mbox{\sc #1}}
      {\mbox{\small\sc #1}}
      {\mbox{\tiny\uppercase{#1}}}%
  \else
    {\sc #1}%
  \fi}



% for call-by-value 
\newcommand{\Lao}{\Lambda} %^{\rm o}}  % all lambda-term
\newcommand{\LaoP}{\Lambda^{+}}  % extended lambda-terms

\newcommand{\LasV}{\Lambda_{\rm V}} % call-by-value lambda-calculus 
\newcommand{\LasN}{\Lambda_{\rm N}} %  call-by-name  lambda-calculus 


\newcommand{\red}{\ensuremath\longrightarrow}

\newcommand{\val}{V}
\newcommand{\valb}{U}
\newcommand{\valp}{V'}

\newcommand{\V}{\mathcal{V}}

\newcommand{\evctxt}{\ensuremath{E  }} %C_{\rm{e}}}} % evaluation contexts in piI
\newcommand{\evctxtp}{\ensuremath{E'} } % _{\rm{e}}}}

% special processes 

\newcommand{\link}[2]{ #1 \vartriangleright #2} % links in call-by-name 
\newcommand{\fwdsymb}{\triangleright}
\newcommand{\fwd}[2]{#1\fwdsymb #2}



% for spacing and similar
\newcommand{\mysp}{15pt}
\newcommand{\mypt}{12pt} 
\newcommand{\myptD}{2pt} 
\newcommand{\myptt}{18pt} 
\newcommand{\myptS}{5pt} 
\newcommand{\myptSmall}{1pt}
\newcommand{\andalso}{\hskip .5cm } 

\newcommand{\reff}[1]{(\ref{#1})} 

% calculi

\newcommand{\Intp}{\mbox{I$\pi$}\xspace}   % Internal pi
\newcommand{\piI}{\Intp}   % Internal pi
\newcommand{\piIop}{\mbox{I$\pi^{\rmm{op}}$}\xspace}   % Internal pi


% substitutions  (used as a postfixed operator)
\def\sub#1#2{\{\raisebox{.5ex}{\small$#1$}\! / \!\mbox{\small$#2$}\}} 

\newcommand{\seb}{\mathrel{\sim}} % strong bisimilarity 


\def\iaeF#1#2{{#1}\,{#2}} % input action in the early style without brackets


%spacial processes for the encoding of $\lambda$-calculus 
\newcommand{\qOUT}{{\mathcal O}} 
\newcommand{\OUT}[5]{\app{\qOUT_{#1}}{#2, #3, #4, \ldots, #5 }}
\newcommand{\OUTU}[5]{\app{\qOUT_{#1}}{#2, #3, #4  }}
\newcommand{\OUTn}[4]{\OUT n{#1}{#2}{#3}{#4 }}
\newcommand{\OUTe}[5]{\OUT{#1}{#2}{#3}{\enco{#4}}{\enco{#5}}}
\newcommand{\OUTen}[4]{\OUT n{#1}{#2}{\enco{#3}}{\enco{#4}}}
\newcommand{\OUTeu}[3]{\app{\qOUT_1}{{#1},{#2},{\enco{#3}} }  }

% END ORIGINAL INPUT davidemacros.sty
% ORIGINAL INPUT guilhemmacros.sty
%% Restored, to be removed
\newcommand{\ired}[1]{\xrightarrow{#1}_i}
\newcommand{\boldemph}[1]{\emph{\textbf{#1}}}
\newcommand{\F}{\mathrm{F}}
\newcommand{\Tr}[1]{\mathbf{Tr}(#1)}
\newcommand{\initName}{}


% taken from the enf section

\newcommand{\enfM}{\mathcal{M}}
\newcommand{\enfK}{\mathcal{K}}
\newcommand{\enfV}{\mathcal{V}}
\newcommand{\enfC}{\mathcal{C}}
\newcommand{\enfG}{\mathcal{G}}
% \newcommand{\enfB}{\mathcal{B}}


\ifnotes
\newcommand{\gj}[1]{\finish{{\bf GJ:} #1}}
\newcommand{\ds}[1]{\finish{{\bf DS:} #1}}
\newcommand{\dsJ}[1]{\finish{{\bf DS(jun23):} #1}}
\newcommand{\gjOLD}[1]{} % comments for now treated. To be
                          %  reconsidered it there is time
\newcommand{\dsOLD}[1]{} % comments for now treated. To be
                          %  reconsidered it there is time


\newcommand{\DSOCT}[1]{\finish{{\bf DS(oct23):} #1}}
\newcommand{\DSOCTb}{\finish{{\bf DS(oct23):} changed below}}
\newcommand{\DSOCTe}{\finish{{\bf DS(oct23):} end changes}}
\else
\newcommand{\gj}[1]{}
\newcommand{\ds}[1]{}
\newcommand{\dsJ}[1]{}
\newcommand{\gjOLD}[1]{} 
\newcommand{\dsOLD}[1]{} 
\newcommand{\DSOCT}[1]{}
\newcommand{\DSOCTb}{}
\newcommand{\DSOCTe}{}
\fi


\newcommand{\N}{\mathbf{N}}
\newcommand{\defeq}{\mathrel{\triangleq}}
% \newcommand{\sep}{~|~}
\newcommand{\dom}[1]{\mathrm{dom}(#1)}
\newcommand{\Power}[1]{\mathcal{P}(#1)}
\newcommand{\emptymap}{\varepsilon}
\newcommand{\emptylist}{\varepsilon}
\newcommand{\parrow}{\rightharpoonup}
\newcommand{\pmap}[2]{[#1 \mapsto #2]}

\newcommand{\Values}{\mathrm{Vals}}
\newcommand{\Terms}{\mathrm{Terms}}
\newcommand{\EContexts}{\mathrm{ECtxs}}
\newcommand{\Contexts}{\mathrm{Ctxs}}

\newcommand{\Var}{\mathrm{Var}}
\newcommand{\FV}{{\tt fv}}
\newcommand{\redv}{\rightarrow_{\tt{v}}}
\newcommand{\Redv}{\Rightarrow_{\tt{v}}}
\newcommand{\redn}{\rightarrow_{\tt{n}}}
\newcommand{\subst}[2]{\{#2 / #1 \}}
\newcommand{\substF}[1]{\{#1\}}

\newcommand{\TLNames}{\mathrm{TLNames}}
\newcommand{\VNames}{\mathrm{VNames}}
\newcommand{\CNames}{\mathrm{CNames}}
\newcommand{\finalName}{\square}
% \newcommand{\init}{\iota}
% \newcommand{\Names}{\mathrm{Names}}
% \newcommand{\lift}[1]{\uparrow\!{#1}}

\newcommand{\act}{\mu}
\newcommand{\Actions}{{\tt Actions}}
% \newcommand{\Init}{{\tt Init}}
\newcommand{\ansP}[2]{\bar{#1}(#2)}
\newcommand{\questP}[2]{\bar{#1}(#2)}
\newcommand{\questPV}[3]{\bar{#1}(#2,#3)}
\newcommand{\questPT}[2]{\bar{#1}(#2)}
\newcommand{\ansO}[2]{#1(#2)}
\newcommand{\questO}[2]{#1(#2)}
\newcommand{\questOV}[3]{#1(#2,#3)}
\newcommand{\questOT}[2]{#1(#2)}
\newcommand{\questOinit}[1]{(#1)}

% \newcommand{\Traces}{{\tt Traces}}
% \newcommand{\JT}{{\tt JT}}
\newcommand{\emptytrace}{\varepsilon}

% \newcommand{\UnivArena}{\mathcal{U}}

% \newcommand \pview[1] {\left\ulcorner #1 \right\urcorner}
% \newcommand \oview[1] {\left\llcorner #1 \right\lrcorner}
\newcommand \sem [1] {\left\llbracket #1 \right\rrbracket}

% \newcommand{\C}{\mathbf{C}}
% \newcommand{\CC}{\mathbf{C}}
% \newcommand{\CCc}{\mathbf{D}}
% \newcommand{\Fc}{\mathrm{G}}
\newcommand{\conf}[1]{\langle #1\rangle}
\newcommand{\initconf}[2]{\langle[? \mapsto #2], #1\rangle}
\newcommand{\mergeConf} [2] {#1 \doublewedge #2}
% \newcommand{\iRed}[1]{\xRightarrow{#1}}
% \newcommand{\sred}[1]{\xrightarrow{#1}_{\mathrm{s}}}
% \newcommand{\sRed}[1]{\xRightarrow{#1}_{\mathrm{s}}}
% \newcommand{\cred}[1]{\xrightarrow{#1}_{\mathrm{c}}}
% \newcommand{\cRed}[1]{\xRightarrow{#1}_{\mathrm{c}}}
% \newcommand{\ared}[1]{\xrightarrow{#1}_{\mathrm{a}}}
% \newcommand{\aRed}[1]{\xRightarrow{#1}_{\mathrm{a}}}
\newcommand{\orderCont}[1]{\prec_{#1}}

\newcommand{\InName}[1]{\widehat{#1}}

% \newcommand{\TrWB}[1]{\mathbf{Tr}_{\tt{wb}}(#1)}

% \newcommand{\iredwb}[1]{\xrightarrow{#1}_{\tt{wb}}}
\newcommand{\singletonStack}[1]{[#1]}
\newcommand{\eraseStack}[1]{\widehat{#1}}

% \newcommand{\comp}[1]{\mathbf{comp}(#1)}

% \newcommand{\restrict}{\upharpoonright}
% \newcommand{\compstrat}[2]{#1;#2}
% \newcommand{\paracompstrat}[2]{#1||#2}
% \newcommand{\idstrat}{{\bf id}}
% \newcommand{\pairingstrat}[2]{\langle#1,#2\rangle}
% \newcommand{\appstrat}{{\bf app}}

% \newcommand{\fred}[1]{\xrightarrow{#1}_f}
% \newcommand{\hred}[1]{\xrightarrow{#1}_h}
% \newcommand{\TrH}[1]{\mathbf{Tr}_h(#1)}
% 
% \newcommand{\TrA}[1]{\mathbf{Tr}_a(#1)}

\newcommand{\lambdarho}{\lambda\rho}
\newcommand{\Loc}{\mathrm{Loc}}
\newcommand{\Stores}{\mathrm{Stores}}
\newcommand{\support}[1]{\mathrm{Sup}_{#1}}
% \newcommand{\ctxapprox}{\lesssim_{ctx}}
\newcommand{\ctxeq}{\simeq_{ctx}}
% \newcommand{\ciuapprox}{\lesssim_{ciu}}
\newcommand{\ciueq}{\simeq_{ciu}}
\newcommand{\ctxeqlr}{=_{\rm{ctx}}^{\lambda \rho}}

\newcommand{\Label}{{\rm pol}}

\newcommand{\cs}{\mathrm{cs}}

\newcommand{\callcc}{\mathrm{call/cc}}

\newcommand{\liftRel}[1]{\widehat{#1}}
\DeclareMathOperator{\Diverge}{\Uparrow}
\newcommand{\Converge}{\mathrel{\Downarrow}}

%%% Pushdown LTS
\newcommand{\pshdsym}{\mathsf{pshd}}
\newcommand{\pshdLTS}{\lts_{\pshdsym}}
\newcommand{\pshdtrans}[1]{\xrightarrow{#1}}%_{\pshdsym}}
\newcommand{\pshdTrans}[1]{\xRightarrow{#1}}%_{\pshdsym}}
\newcommand{\cnstack}{\xi}
\newcommand{\cnstackb}{\xi'}
\newcommand{\emptystack}{[]}
\newcommand{\cnset}{\Sigma}
% END ORIGINAL INPUT guilhemmacros.sty

\theoremstyle{plain}
\newtheorem{theorem}[thm]{Theorem}
\newtheorem{proposition}[thm]{Proposition}
\newtheorem{lemma}[thm]{Lemma}
\newtheorem{corollary}[thm]{Corollary}
\newtheorem{conjecture}[thm]{Conjecture}
\theoremstyle{definition}
\newtheorem{definition}[thm]{Definition}
\newtheorem{example}[thm]{Example}
\newtheorem{remark}[thm]{Remark}

\ifnotes
\newcommand{\DS}[1]{[ DS: \emph{#1}]}
\else
\newcommand{\DS}[1]{}
\fi

\ifcommentreview
\newcommand{\reviewComment}[1]{\textbf{Review}: {\color{red} #1}}
\newcommand{\reviewCommentDS}[1]{\textbf{Review[DS]}: {\color{green} #1}}
\newcommand{\reviewCommentGJ}[1]{\textbf{Review[GJ]}: {\color{blue} #1}}
\else
\newcommand{\reviewComment}[1]{}
\newcommand{\reviewCommentDS}[1]{}
\newcommand{\reviewCommentGJ}[1]{}
\fi

\newcommand{\cutLMCSsecondRound}[1]{}

%I agree about removing the full abstraction for lambda store from the
%paper, including Section 13. So I have just done so, removing all the
%references about the calculus and the full-abstraction result from
%the paper (using the \cutLMCSsecondRound macro). 



\def\doublewedge{\mathrel{\wedge\!\!\!\wedge}}




% \citestyle{acmauthoryear}
%% Copyright information
%% Supplied to authors (based on authors' rights management selection;
%% see authors.acm.org) by publisher for camera-ready submission;
%% use 'none' for review submission.
%\setcopyright{none}
%\setcopyright{acmcopyright}
%\setcopyright{acmlicensed}
%\setcopyright{rightsretained}
%\copyrightyear{2017}           %% If different from \acmYear

%% Bibliography style
%\bibliographystyle{ACM-Reference-Format}
%% Citation style
%\citestyle{acmauthoryear}  %% For author/year citations
%\citestyle{acmnumeric}     %% For numeric citations
%\setcitestyle{nosort}      %% With 'acmnumeric', to disable automatic
                            %% sorting of references within a single citation;
                            %% e.g., \cite{Smith99,Carpenter05,Baker12}
                            %% rendered as [14,5,2] rather than [2,5,14].
%\setcitesyle{nocompress}   %% With 'acmnumeric', to disable automatic
                            %% compression of sequential references within a
                            %% single citation;
                            %% e.g., \cite{Baker12,Baker14,Baker16}
                            %% rendered as [2,3,4] rather than [2-4].


%%%%%%%%%%%%%%%%%%%%%%%%%%%%%%%%%%%%%%%%%%%%%%%%%%%%%%%%%%%%%%%%%%%%%%
%% Note: Authors migrating a paper from traditional SIGPLAN
%% proceedings format to PACMPL format must update the
%% '\documentclass' and topmatter commands above; see
%% 'acmart-pacmpl-template.tex'.
%%%%%%%%%%%%%%%%%%%%%%%%%%%%%%%%%%%%%%%%%%%%%%%%%%%%%%%%%%%%%%%%%%%%%%
%\usepackage[T1]{fontenc}
%\usepackage{lmodern}


%% Some recommended packages.
\usepackage{booktabs}   %% For formal tables:
                        %% http://ctan.org/pkg/booktabs
\usepackage{subcaption} %% For complex figures with subfigures/subcaptions
                        %% http://ctan.org/pkg/subcaption




\AtBeginDocument{\def\shorttitle{Games, processes and functions}}
